
\subsubsection{Temperature Scaling and Post-Hoc Calibration}

Temperature scaling, introduced by Guo et al. (2017), optimizes a single parameter T to rescale logits before softmax, improving Expected Calibration Error (ECE) for neural networks. This method incurs zero inference overhead since T is applied post-training. Guo et al. focused on calibration metrics (ECE), not selective prediction AUROC. Nixon et al. (2019) explored ensemble temperature scaling, and Desai \& Durrett (2020) studied calibration in natural language generation tasks. These works do not report AUROC versus cost trade-offs, treating calibration and selective prediction as separate problems.

\subsubsection{Conformal Prediction for Language Models}

Conformal prediction provides distribution-free coverage guarantees for uncertainty sets. Kumar et al. (2023) introduced conformal prediction for selective natural language generation, demonstrating validity under distribution shift. Su et al. (2024) extended this to API-only settings using sample frequency and semantic similarity for nonconformity scores, achieving performance on hallucination detection. These methods establish conformal prediction viability for LLMs but do not compare inference cost across multiple UQ approaches.

\subsubsection{Monte Carlo Dropout}

Gal \& Ghahramani (2016) showed that dropout at inference time approximates Bayesian inference, enabling uncertainty estimation via variance across k stochastic forward passes. This approach has been applied to computer vision \citep{Kendall2017Uncertainties} and natural language processing \citep{Xiao2019Dropout}, typically using k=10-30 samples. Recent work on LLM uncertainty includes Kuhn et al. (2023) on semantic entropy and Lin et al. (2023) on conformal language modeling. However, k-dependency for LLMs remains understudied. Prior work does not systematically vary k to identify efficiency points.

\subsubsection{Selective Prediction Benchmarks}

Yang et al. (2023) applied selective prediction to TruthfulQA, showing that uncertainty-based rejection improves accuracy. Their work validates TruthfulQA as a selective prediction benchmark but does not compare multiple UQ methods or report inference costs. Zhang et al. (2020) compared ensemble and compositional calibration methods, introducing kernel-density ECE estimators. This work predates LLM-scale models and lacks selective prediction evaluation.

\subsubsection{Our Positioning}

Existing work evaluates UQ methods in isolation or compares methods without cost analysis. We are the first to systematically benchmark AUROC versus inference cost for multiple UQ method variants on the same selective prediction task. Our Pareto frontier framing shifts the question from ''which method wins'' to ''which method for my budget.''
